\section{Removing Digits}
\label{sec:cses-removing-digits}

\problemstatement{Start from an integer $n$. On each step you may
subtract any one of the digits currently appearing in the number. Find
the minimum number of steps to reduce $n$ to $0$.}

\begin{patternbox}
\emph{``Minimum steps to reduce a number to zero''} is a shortest-path
DP over values $0\ldots n$. Let $\textit{dp}[v]$ be the fewest steps from
$v$ to $0$; from $v$ you may subtract any digit of $v$, so
$\textit{dp}[v]=1+\min_{d\in\text{digits}(v)}\textit{dp}[v-d]$.
\end{patternbox}

\subsection*{State and transition}

$\textit{dp}[v]$ = minimum steps to bring $v$ down to $0$, with
$\textit{dp}[0]=0$. The only moves from $v$ subtract one of $v$'s own
decimal digits. Since a larger digit removes more in one step but the
available digits depend on $v$ itself, we must try every digit of $v$ and
take the minimum successor:
\[
  \textit{dp}[v]=1+\min_{d\,\in\,\text{digits}(v),\ d>0}\textit{dp}[v-d].
\]

\begin{ideabox}[Greedy ``Subtract the Largest Digit'' Also Works]
Because subtracting a larger digit never overshoots ($v-d\ge0$ always,
as $d\le9\le v$ for $v\ge10$, and single digits vanish in one step), a
greedy that always removes the largest digit is optimal here. The DP is
shown because it generalises and needs no proof of the greedy exchange
argument.
\end{ideabox}

\subsection*{Algorithm}

\begin{enumerate}
  \item $\textit{dp}[0]=0$.
  \item For $v=1\ldots n$: for each non-zero digit $d$ of $v$,
        $\textit{dp}[v]=\min(\textit{dp}[v], \textit{dp}[v-d]+1)$.
  \item Answer $\textit{dp}[n]$.
\end{enumerate}

\subsection*{Dry run}

\dryrun{$n=27$.}

\begin{itemize}
  \item From $27$ subtract $7\to20$; from $20$ subtract $2\to18$; then
        $8\to10$; then $1\to9$; then $9\to0$.
  \item The DP finds $\textit{dp}[27]=5$ via the value chain
        $27\to20\to18\to10\to9\to0$ (each step removes one digit).
\end{itemize}

\subsection*{C++ implementation}

\begin{lstlisting}
#include <bits/stdc++.h>
using namespace std;

int main() {
    int n;
    cin >> n;
    vector<int> dp(n + 1, 1e9);
    dp[0] = 0;
    for (int v = 1; v <= n; v++) {
        int t = v;
        while (t > 0) {
            int d = t % 10;
            if (d > 0) dp[v] = min(dp[v], dp[v - d] + 1);
            t /= 10;
        }
    }
    cout << dp[n] << "\n";
    return 0;
}
\end{lstlisting}

\begin{complexitybox}
\textbf{Time Complexity.} $O(n\log_{10} n)$ -- each value inspects its
$O(\log n)$ digits. \textbf{Space Complexity.} $O(n)$ for the DP array.
\end{complexitybox}

\begin{takeawaybox}
When moves from a state depend on the state's own structure (here, its
digits), enumerate those moves inside the transition. This ``value as a
node, allowed subtractions as edges'' framing turns a reduce-to-zero
puzzle into a one-dimensional shortest-path DP.
\end{takeawaybox}
